\begin{frame}
  \frametitle{课程安排 | 授课教材}
  \begin{figure}
    \centering
    \includegraphics[width=6cm]{c1.intro.textbook.png}
  \end{figure}
\end{frame}

\begin{frame}
  \frametitle{课程安排 | 理论课}
  \begin{center}
  \alert{前9周，每周一\&五，上午前两节（8:00-9:50），西楼610}\\
  \vspace{0.2cm}
  \end{center}
  \begin{block}{授课内容}
    \begin{itemize}
      \item 教材内容：第一章……～第十章
      \item 补充知识：Markdown, Git, ...
    \end{itemize}
  \end{block}
\end{frame}

\begin{frame}
  \frametitle{课程安排 | 理论课 | 自主学习}
  \begin{itemize}
    \item 上课期间以小组形式编写程序
    \item 任务主题不限（生信方面的最好）
    \item 理论课的最后一/两次课（2～4学时）进行报告
    \item 小组成员人人都要参与
    \item 既参与编程又进行报告者加分
    \item 作为平时成绩的一大部分
  \end{itemize}
\end{frame}

\begin{frame}
  \frametitle{课程安排 | 理论课 | 自主学习 | 实例}
  \begin{block}{目的}
    按照用户的需求生成随机密码。
  \end{block}
  \pause
  \begin{block}{需求}
    \begin{itemize}
      \item 指定生成密码的元素（默认：混合使用数字、大小写字母、符号）
      \item 指定生成密码的长度（默认：12个字符）
      \item 指定生成密码的个数（默认：1个）
      \item 指定密码中元素的比例/字符数（默认：无比例/字符数要求）
      \item 指定需要排除的元素（比如：0和O，1和l；默认：无需排除的元素）
      \item 挖掘更多人性化需求……
    \end{itemize}
  \end{block}
\end{frame}

\begin{frame}
  \frametitle{课程安排 | 实验课}
  \begin{center}
  \alert{前9周，每周一，上午后两节（10:10-12:00），教一楼304}\\
  \vspace{0.2cm}
  \end{center}
  \begin{block}{实验内容}
    \begin{itemize}
      \item 紧跟理论课进度
      \item 教材上的程序/实例
    \end{itemize}
  \end{block}
\end{frame}

\begin{frame}
  \frametitle{课程安排 | 考核方式}
  \begin{enumerate}
    \item 理论课：60\%
      \begin{enumerate}
        \item 平时表现：10\%
        \item 闭卷考试：50\%
      \end{enumerate}
    \item 实验课：40\%
      \begin{enumerate}
        \item 平时表现：20\%
        \item 实验报告：20\%
      \end{enumerate}
  \end{enumerate}
\end{frame}

